\section{Same seed, same flower}\label{sec:determinism}
\stagebox{\statusbuilt}{\src{core/src/math.rs}, \src{core/src/rng.rs}, \src{checks/}, \src{.github/workflows/ci.yml}}

A seed must give the same bud, bit for bit, on every machine and in every
engine, including the browser build. Three rules make that true.

\begin{enumerate}[leftmargin=*]
\item \textbf{One specified generator.} All randomness comes from the PCG32
  streams of \cref{sec:seeds}, never from a library or the operating system.
\item \textbf{Portable maths.} Platform maths libraries can differ in the last
  bit of $\exp$, $\ln$, $\sin$, $\cos$, $x^y$ and $\operatorname{atan2}$, and
  those differences grow through thousands of relaxation steps. All
  transcendental functions go through the \texttt{libm} crate, a pure-Rust
  implementation that gives the same bits everywhere. Square roots and basic
  arithmetic are correctly rounded by IEEE~754 on every target, so square
  roots use the hardware instruction (\texttt{libm}'s own falls back to a slow
  software routine under WebAssembly).
\item \textbf{No hidden ordering.} Floating-point addition is not associative,
  so nothing that reaches the output may depend on hash-map iteration order or
  on threads summing in varying order. The contact search sorts its result,
  so its spatial grid only affects speed (\cref{sec:packing}).
\item \textbf{Threads that cannot change the bits.} Contact search and
  contact detection run on all cores natively (\src{core/src/par.rs}). The
  work is cut into a fixed number of chunks that depends only on its size,
  never on the core count, and the chunks' results are joined in chunk order.
  Every floating-point sum then runs in one thread in that order, so a flower
  has the same bits on one core or sixteen. WebAssembly runs the same chunks
  one after another.
\end{enumerate}

\subsection{How it is checked}
\begin{itemize}[leftmargin=*]
\item The canonical output for a seed is the flower JSON of flower 0, with
  the peony preset and default settings (\texttt{grow\_bud\_json(seed, 0)}).
  The check grows its seeds on parallel threads, and each bud may use
  threads of its own (see above); neither can change its bits.
\item \src{checks/golden/digests.txt} holds a SHA-256 of that output for eight
  reference seeds, chosen to cover small, large and bit-pattern values: 0, 1,
  2, 7, 42, 1234, \texttt{0xDEADBEEF} and $2^{64}-1$.
\item CI computes the digests on Linux, macOS, Windows and WebAssembly and
  fails if any differ from the file.
\item A change that is meant to alter flowers re-runs the check with
  \texttt{BLESS=1} and commits the new file with the change, so the diff shows
  which seeds changed.
\end{itemize}
